%==============================================================================
\section{From equations to array operations}\label{sec:discretisation}
%==============================================================================

K24 has no linear system to solve. A solve is a fixed sequence of array operations (potential, filling, gradients,
smoothing, routing, conversion to $q$, conduit balance), wrapped in at most one Picard loop. This section describes
each operation as the code performs it.

\subsection{The grid}\label{disc:grid}
\code{ArrayHydroGrid(Nx, Ny, xlims, ylims; T = Float64)} is a cell-centred grid of \code{Nx}\,$\times$\,\code{Ny} cells.
\code{xlims}/\code{ylims} are the outer \emph{edges} of the domain, so $\Dx = (x_2 - x_1)/N_x$, and the centre of
cell $(i,j)$ lies at $x_1 + (i-\tfrac12)\Dx$. The grid stores only \code{Nx}, \code{Ny}, \code{dx}, \code{dy} and an FFT
cache; it has no coordinates, so a plot's axes are measured from the lower-left corner. \code{compute\_lims(xc, yc)}
turns cell-centre coordinates from a data file into these edge limits. Every field is a plain
\code{Matrix\{T\}} of size $(N_x, N_y)$: no halos, no ghost cells. $\Dx\ne\Dy$ is supported everywhere; the code
never collapses the two spacings into one.

The grounded mask is a float field. A cell is grounded if and only if \code{mask[i, j] == 1}; any other value
(0, or a fractional value) is ``not grounded''. The hydrology lives on grounded cells: $N$, $q$, $H$, $\Sinf$, $\Ninf$
are 0 elsewhere, and routing never passes water through a non-grounded cell.

\subsection{Gradients}\label{disc:gradients}
\code{minus\_gradient\_x!} writes $-\partial f/\partial x$ with centred differences inside and one-sided differences on
the two edge columns:
\begin{equation}
-\partial_x f\big|_{i} \approx \begin{cases} -(f_{2}-f_{1})/\Dx & i = 1,\\ -(f_{i+1}-f_{i-1})/(2\Dx) & 1<i<N_x,\\
-(f_{N_x}-f_{N_x-1})/\Dx & i = N_x,\end{cases}
\end{equation}
and 0 if $N_x = 1$; the same in $y$. The edge rule gives the full gradient; the source records that the earlier
rule (replicate the edge cell as a ghost and keep $2\Dx$) halved it. A linear field gets its exact gradient in every
cell, and centred differences are exact for a quadratic inside (tour PART~1). \code{update\_potential\_gradients!}
uses this twice: on the true $\phi_0$ to form $|\nabla\phi_0|$ (Euclidean), which enters $\Sinf$, $\Ninf$, the
dissipation and $W$; and on the filled potential $\phif$ to get the flow direction.

\subsection{Filling the pits of the routing potential}\label{disc:filling}
A cell lower than all four neighbours would trap water. \code{potential\_filling!} copies the routing surface into
\code{phi0\_filled} (by default $\phi_0$; for the non-GDS schemes the smoothed potential) and removes pits by one of
three algorithms (\code{fill\_algorithm}):
\begin{itemize}[leftmargin=1.5em]
\item \code{JacobiFill} (default for the GDS schemes, the original K24/KORI filling): up to \code{fill\_iters} = 10
  Jacobi passes, each raising every strict local minimum to the mean of its four neighbours. Domain edges are
  treated as replicated neighbours, so an edge cell is never a minimum. The mask is ignored. This converges slowly
  and does not guarantee a pit-free surface: on the tour's 80\,m-bump fixture (PART~14) 3 of its 4 pits are still strict minima
  after the 10 passes, where \code{LowestNeighbourFill} removes all of them.
\item \code{LowestNeighbourFill} (Le Brocq Sec.~3.1): the same passes, raising each minimum to its lowest neighbour.
  Leaves flats where hollows were.
\item \code{PriorityFloodFill(; epsilon = 1 Pa)} (Barnes et al.\ 2014; default for the other schemes): the outlets are
  every non-grounded cell next to a grounded one (grounding line, ice margin) and every grounded cell on the domain
  edge. Starting from the lowest outlet, cells are taken in increasing order of filled potential from a binary
  min-heap; each newly reached grounded neighbour not higher than the cell it was reached from is raised to that
  value plus $\epsilon$. Afterwards every grounded cell has a strictly decreasing path to an outlet, in one
  $O(n\log n)$ pass. $\epsilon = 1$\,Pa is $\sim0.1$\,mm of water head. Only grounded cells change.
\end{itemize}
The Jacobi passes are implemented with a candidate list: after the first full-grid pass only the cells filled in
the previous pass and their neighbours can be minima, so later passes test only those
(\code{collect\_fill\_candidates!}, deduplicated with a per-cell pass stamp). The tour checks that this gives the same
surface, bit for bit, as naive full-grid passes. The filled surface is used for the flow direction only; $\phi_0$ is
never modified, so filling does not change $N$. \code{model.h} receives the thickness consistent with $\phif$,
used only for the mean thickness of the smoothing kernel.

\subsection{Stress-gradient smoothing}\label{disc:smoothing}
The smoothing kernel (\code{coupling\_kernel}) is a cone of radius $w$,
\begin{equation}
k(r) = \max\bigl(0,\ 1 - r/w\bigr), \qquad w = 4\,\bar h\,\lambda, \qquad \textstyle\sum k = 1,
\end{equation}
where $r$ is the physical distance from the kernel centre (computed with $\Dx$ and $\Dy$ separately, so the support is
a circle in physical space and the array is rectangular when $\Dx\ne\Dy$), $\bar h = \max(\text{mean grounded
thickness}, 10\,\text{m})$, and $\lambda$ = \code{longcoupwater}, KORI-ULB's parameter. The array has
$2\,\mathrm{round}(w/\Dx - \tfrac12) + 1$ columns in $x$ (likewise in $y$). The area-weighted mean distance of the cone,
$\int_0^w r\,k(r)\,r\,dr / \int_0^w k(r)\,r\,dr = w/2 = 2\bar h\lambda$, is the effective coupling length to compare
with Kamb and Echelmeyer's range (4--10 ice thicknesses for ice sheets, 1--3 for valley glaciers, about 12 in a
surge); a 1-D mean, $w/3$, would understate it. The constructor therefore takes the coupling length directly,
\code{coupling\_length\_kamb86} $= 2\lambda$, so the effective length is \code{coupling\_length\_kamb86}\,$\times\,\bar h$.

How the kernel is applied depends on the scheme (\code{prepare\_routing!}):
\begin{itemize}[leftmargin=1.5em]
\item GDS schemes: fill $\phi_0$, take $-\nabla\phif$, convolve each component with $k$
  (\code{update\_smoothed\_potential\_gradients!}) to get the smoothed driving direction $\mathbf s$ (\code{minus\_grad\_phi0\_sx/sy}).
\item Warner, Quinn, Tarboton, ModifiedTarboton: convolve the potential, then fill the result, then take its gradient.
  Smoothing commutes with differentiation, so this is the same coupling, but the routing follows one surface, so
  downhill-only schemes stay acyclic.
\end{itemize}
With $\lambda = 0$ nothing is convolved. The convolution uses replicate padding at the domain edge, through a cached FFT
(Section~\ref{ref:fft}). The ``smoothed gradient magnitude'' \code{abs\_grad\_phi0\_s} is computed as $|s_x| + |s_y|$, an
$L^1$ norm (it is used only by \code{LaminarThickness}; see Section~\ref{pit:open}).

\subsection{Routing schemes}\label{disc:schemes}
A scheme decides the fractions $w_{n\to c}$ of Eq.~\eqref{eq:psi}. Directions are numbered as in \code{ROUTE\_OFFSETS}:
1 $(-1,0)$ west, 2 $(+1,0)$ east, 3 $(0,-1)$ south, 4 $(0,+1)$ north, 5 SW, 6 SE, 7 NW, 8 NE; \code{ROUTE\_OPPOSITE} maps
each to its reverse. The options are those compared by Le Brocq et al.\ (2006, Table~1):
\begin{center}\small
\begin{longtable}{>{\raggedright\arraybackslash}p{2.6cm}>{\raggedright\arraybackslash}p{1.1cm}>{\raggedright\arraybackslash}p{11.3cm}}
\toprule
scheme & nbrs & outflow fraction from cell $c$ to neighbour $d$ \\
\midrule
\endhead
\code{GDSWarner} (original K24) & 4 & $w_d = \max\bigl(0,\ (s_x d_i\Dy + s_y d_j\Dx)/(|s_x|\Dy + |s_y|\Dx)\bigr)$ with $\mathbf s$ the smoothed
  gradient of $\phif$: at most two neighbours, one per component. The denominator is $c_f|\mathbf s|$, so these fractions are
  consistent with Eq.~\eqref{eq:corfac}. Can send water uphill on $\phif$. \\
\code{Warner} (default; Budd \& Warner 1996, Le Brocq Eq.~8) & 4 & $w_d \propto \max(0,\tilde\phi_{0,c} - \tilde\phi_{0,d})\times(\text{face length}/\text{distance})$,
  i.e.\ $\Dy/\Dx$ for $x$-neighbours and $\Dx/\Dy$ for $y$-neighbours (the face-normal gradient times the face length), normalised to sum to 1. \\
\code{Quinn()} & 8 & $w_d \propto \max(0,\tilde\phi_{0,c}-\tilde\phi_{0,d})$ over 8 neighbours (Le Brocq Eq.~8 as written: no distance factor). \\
\code{Quinn(original = true)} & 8 & Quinn et al.\ (1991): $w_d\propto(\text{drop}/\delta_d)\,L_d$ with contour lengths $L = \Dy/2$, $\Dx/2$ (cardinal)
  and $\tfrac{\sqrt2}{4}\sqrt{\Dx\Dy}$ (diagonal). Equal drops on a square grid give 1/6 per cardinal, 1/12 per diagonal. \\
\code{Tarboton} (D$\infty$, 1997) & 8 & steepest descent over the 8 triangular facets (\code{TARBOTON\_FACETS}, each a cardinal and an adjacent
  diagonal neighbour); the outflow is split between the facet's two neighbours in proportion to the angle. \\
\code{Modified\allowbreak{}Tarboton} & 8 & one direction from the local 4-neighbour gradient of $\phif$, split between the two of 8 neighbours bracketing it (\code{angle\_split!}). \\
\code{GDSTarboton} & 8 & as ModifiedTarboton, with the smoothed gradient components. \\
\bottomrule
\end{longtable}
\end{center}
Water sent to a non-grounded neighbour leaves the system there (the grounding line). Warner and Quinn send water
only downhill on $\phif$, so their routing graph has no cycles, and after \code{PriorityFloodFill} every drop of the
source leaves the grounded domain: routing conserves water exactly (tour PART~15 checks this to $10^{-12}$). The GDS
schemes can send water uphill on $\phif$ (Le Brocq Sec.~3, step~5), which creates cycles; the source records that on the
Thwaites 2\,km, Greenland 8\,km and Antarctica 16\,km test data GDSWarner loses 3\,\%, 8\,\% and 34\,\% of the positive
basal melt in them, Warner none. This is why Warner became the default.

Every scheme except GDSWarner stores its fractions in \code{routing\_tape.w8[d, i, j]}; GDSWarner computes them on the
fly from \code{routing\_weight}. The options that follow a scheme by default (left as \code{nothing} in the constructor):
\begin{center}\small
\begin{tabular}{>{\raggedright\arraybackslash}p{4.0cm}>{\raggedright\arraybackslash}p{3.2cm}>{\raggedright\arraybackslash}p{3.2cm}>{\raggedright\arraybackslash}p{4.2cm}}
\toprule
scheme & \code{fill\_algorithm} & \code{q\_conversion} & \code{dissipation\_discretization} \\
\midrule
\code{Warner} & \code{Priority\allowbreak{}Flood\allowbreak{}Fill} & \code{QFrom\allowbreak{}Face\allowbreak{}Average} & \code{FaceDissipation} \\
\code{GDSWarner}, \code{GDSTarboton} & \code{JacobiFill} & \code{QFromOutflow} & \code{Cell\allowbreak{}Centred\allowbreak{}Dissipation} \\
\code{Quinn}, \code{Tarboton}, \code{Modified\allowbreak{}Tarboton} & \code{Priority\allowbreak{}Flood\allowbreak{}Fill} & \code{QFromOutflow} & \code{Cell\allowbreak{}Centred\allowbreak{}Dissipation} \\
\bottomrule
\end{tabular}
\end{center}
So \code{routing\_scheme = GDSWarner()} alone reproduces the original K24 set-up.

\subsection{The routed flux}\label{disc:psi}
Eq.~\eqref{eq:psi} defines $\psiout$ recursively: a cell's outflow needs its upstream neighbours' outflows first.
The reference implementation (\code{accumulate\_psi\_out!}, KORI-ULB's \code{DpareaWarGds.m}) is a depth-first
recursion with a \code{visited} flag; Section~\ref{sec:routing} describes it and its three alternatives. Three rules
complete the definition:
\begin{itemize}[leftmargin=1.5em]
\item A neighbour off the domain contributes nothing (the no-flux divide, K24 Eq.~2b).
\item $\psiout\ge0$: a cell whose melt is negative enough to refreeze all its inflow passes on nothing.
\item At most \code{max\_psi\_out\_calls} = $10^5$ cells are visited per sweep (KORI-ULB caps at $5\times10^4$); past the cap
  a cell keeps only its own source. This is a guard against runaway traversal, and is large enough for the test
  grids (Thwaites 2\,km has about 52\,000 grounded cells).
\end{itemize}

\subsection{From routed flux to $q$}\label{disc:qconv}
Two conversions (\code{q\_conversion}):
\begin{itemize}[leftmargin=1.5em]
\item \code{QFromOutflow}: $q = \psiout/(c_f + \varepsilon)$, Eq.~\eqref{eq:corfac}, clamped to $[q_{\min}, q_{\max}]$. It
  describes the flux at the cell's downstream side. $\varepsilon$ = \code{eps(T)}: where the smoothed gradient vanishes,
  $c_f = 0$ and the bare division gave NaN or Inf.
\item \code{QFromFaceAverage} (4-neighbour schemes only; default with Warner): the face fluxes
  \begin{equation}
  F^x_{a,j} = \psiout_{a-1,j}\,w_{2,(a-1,j)} - \psiout_{a,j}\,w_{1,(a,j)},\qquad
  F^y_{i,b} = \psiout_{i,b-1}\,w_{4,(i,b-1)} - \psiout_{i,b}\,w_{3,(i,b)},
  \end{equation}
  (net, positive in $+x$/$+y$, on staggered arrays of size $(N_x{+}1)\times N_y$ and $N_x\times(N_y{+}1)$; only grounded
  cells contribute) give face-normal fluxes $F^x/\Dy$, $F^y/\Dx$, averaged to the centre component by component:
  $q = \bigl|\bigl(\tfrac{F^x_{i}+F^x_{i+1}}{2\Dy},\ \tfrac{F^y_{j}+F^y_{j+1}}{2\Dx}\bigr)\bigr|$, the mean of inflow and outflow.
\end{itemize}
With face fluxes, the net outflow through a cell's four faces equals its source, $\mdot\Dx\Dy/\rhow$, wherever its
outflow fractions sum to 1: a discrete divergence theorem (tour PART~15).

\subsection{The dissipation melt}\label{disc:dissipation}
\begin{itemize}[leftmargin=1.5em]
\item \code{CellCentredDissipation}: $Q_{\rm diss} = |q|\,|\nabla\phi_0|$ at the cell centre, from the true potential (not
  the smoothed one: smoothing is a device for the flow direction, not the local work done on the water).
\item \code{FaceDissipation} (4-neighbour schemes; default with Warner): on each face the water releases
  $F(\phi_{\rm up}-\phi_{\rm down})$ [W] of the true potential; each grounded cell receives half the energy of each of its
  in-domain faces: $Q_{\rm diss} = \tfrac{1}{2\Dx\Dy}\sum_{\rm faces} F\,\Delta\phi$. The sum over cells equals the energy
  released on all faces (a face on the grounding line gives its grounded cell half). It is signed: water pushed up
  the true potential (out of a filled hollow) gives a negative term.
\end{itemize}

\subsection{Frictional heat}\label{disc:friction}
\code{CellCentredFriction} (default): $Q_b = \tau_b\,|u_b|$ at the centre. \code{StaggeredFriction(ux, uy; quadrature, u\_floor)}
takes C-grid velocities as Yelmo provides them (\code{ux[i,j]} on the face between cells $i$ and $i{+}1$, Yelmo's
\code{acx}; \code{uy} likewise, \code{acy}). The sliding law is still evaluated at the centre; its drag
$\beta = \tau_b/\max(|u_b|, u_{\rm floor})$ is averaged to each face, giving face tractions $\tau_x = \beta_{\rm face}u_x$,
$\tau_y = \beta_{\rm face}u_y$ (the way Yelmo builds \code{taub\_acx/acy}). Then either
\begin{equation}
Q_b = \tfrac12\bigl(\tau_x u_x|_{i-1/2} + \tau_x u_x|_{i+1/2}\bigr) + \tfrac12\bigl(\tau_y u_y|_{j-1/2} + \tau_y u_y|_{j+1/2}\bigr)
\quad(\text{faces, } \beta u^2\ge0),
\end{equation}
or, with \code{quadrature = true} (Yelmo's \code{qb\_method = 2}), the face velocities and tractions are interpolated
bilinearly to the four points of a $2\times2$ Gauss rule in the cell and $|\mathbf u||\boldsymbol\tau|$ is averaged there.
Edge faces use the one cell they have; $u_{\rm floor} = 10^{-3}$\,m/yr (Yelmo's \code{ub\_sq\_min}). For a uniform flow
both give $\tau_b|u_b|$ exactly (tour PART~13).

\subsection{The effective-pressure kernel}\label{disc:N}
\code{update\_N!} evaluates Eqs.~\eqref{eq:Sinf}--\eqref{eq:N} cell by cell in one pass over grounded cells
(\code{update\_N\_grounded\_column!}, columns split across threads). Four degenerate cases are resolved explicitly, each
in favour of its physical limit:
\begin{itemize}[leftmargin=1.5em,itemsep=0pt]
\item $Q = 0$ (and $|\nabla\phi_0| = 0$): $\Sinf = 0$, not $0^{-0.4}\cdot0^{0.8} = \infty\cdot0$ = NaN.
\item $Q_c = 0$ (\code{EfficientOnly}) and $Q = 0$: $H_{\rm soft} = 0$, not $e^{-0/0}$ = NaN.
\item $\Sinf = 0$: $\Ninf = p_o$ (no conduit, no drainage).
\item $\Ninf = 0$: $N = 0$, not $\operatorname{erf}(\phi_0/0)\cdot0$.
\end{itemize}
Non-grounded cells get $N = 0$ and zero conduit fields; $p_o = \rhoi g h$ is written everywhere.

\subsection{Freeze-on capacity}\label{disc:freeze}
For an ice model whose basal boundary condition can refreeze water, \code{freeze\_on\_capacity!} returns how much water
reaches each grounded cell [m/s of ice]:
\begin{equation}
C = \frac{1}{\rhoi}\max\Bigl(0,\ \frac{\rhow\,\Psi_{\rm in}}{\Dx\Dy} + \ieb\Bigr),\qquad
\Psi_{\rm in}(c) = \sum_{n} \psiout(n)\,w_{n\to c},
\end{equation}
the inflow from upstream neighbours plus the water from above. The cell's own melt and dissipation are excluded:
they are heat, which the host already counts in its freezing demand.
